% The content of the blueprint, shared by the web and the print version.
%
% Every entry gives the statement in LaTeX next to the exact Lean name (\lean{...}), which
% `lake exe checkdecls blueprint/lean_decls` checks against the compiled project. \leanok marks a
% statement that is formalized and a proof that is complete (the project has no `sorry`).
% Named hypotheses (the arrow, the decision hypothesis, law 2.2) are definitions of type `Prop`:
% \leanok on them means that the definition exists, never that the hypothesis is proved.

\chapter*{Introduction}

This blueprint accompanies the Lean 4 artifact of the essay \emph{The Computational Arrow of Time:
Locating the Second Law of Computation} (in preparation). It puts every definition and theorem of the
chain
\[
  \text{decisions} \to \text{laws} \to \text{localization} \to \text{second law} \to
  \text{consequences for } \cNP
\]
next to its name in Lean, so that a reader can check that the theorem verified is the theorem
stated, and not one that only looks like it.

Two layers carry the argument. The \emph{core} (\texttt{Sagitta.Arrow}, library
\texttt{SagittaCore}, only Mathlib) proves it for every abstract model with a few closure
properties. The \emph{Turing-machine layer} (\texttt{Sagitta.Arrow.TM}, library
\texttt{SagittaTM}) shows that the multi-tape machines of cslib, with the classes of
cslib-frontier, form such a model, proves Bennett's theorem for them, and derives the
consequences with the vendored proofs of Cook--Levin, Ladner and search-to-decision.

\textbf{What a green node means.} Every node of the dependency graph is formalized, and every proof
is complete: the project has no \texttt{sorry}, and each main file ends with
\texttt{\#print axioms} under \texttt{\#guard\_msgs}, so only \texttt{propext},
\texttt{Classical.choice} and \texttt{Quot.sound} are used. The three named hypotheses --- the
arrow ($\cP \neq \cNP$), the decision hypothesis and law 2.2 --- are \emph{definitions} of type
\texttt{Prop}, taken as explicit arguments; a green node for one of them says that the definition
exists, not that the hypothesis holds. None of them is a Lean \texttt{axiom}.

\textbf{Sources.} An entry that formalizes a result of the literature names its source; when the
Lean statement differs from the source (another machine model, another encoding, a weaker or a
stronger form), a \emph{Fidelity} note says how. The laws, the localization and the decision
hypothesis are this project's own; they are marked as such and have no source.

\chapter{The model and the classes}
\label{chap:model}

The argument is stated once for an abstract model and once for Turing machines. The abstract
model only names the classes and the closure properties that the argument uses; the
Turing-machine instance takes them from cslib and cslib-frontier, with their proofs.

\section{The abstract model}

\begin{definition}[Abstract complexity model]
  \label{def:model}
  \lean{Sagitta.Arrow.Model}
  \leanok
  A \emph{model} consists of a type of decision problems and a type of search problems; predicates
  $\cP \subseteq \cNP$ on decision problems and $\cFP$, $\cFNP$ on search problems; a reflexive and
  transitive relation $A \red B$ of reductions; and two maps $\dom$ and $\pre$ from search problems
  to decision problems (the domain, ``does the instance have a solution?'', and the prefix problem,
  ``does some solution of $a$ begin with $u$?''). They satisfy:
  \begin{enumerate}
    \item if $A \red B$ and $B \in \cP$ then $A \in \cP$ (an efficient procedure may call an
      efficient subroutine);
    \item if $R \in \cFNP$ then $\dom R \in \cNP$ and $\pre R \in \cNP$; if $R \in \cFP$ then
      $\dom R \in \cP$;
    \item every $A \in \cNP$ is $\dom R$ for some $R \in \cFNP$;
    \item if $R \in \cFNP$ and $\pre R \in \cP$ then $R \in \cFP$ (search through prefixes).
  \end{enumerate}
  \own{} The closure properties are standard facts about the classes of the literature: see
  \cite{AroraBarak} (Chapter~2) and \cite{Goldreich} (Chapter~2, search and decision).
\end{definition}

\begin{definition}[$\cP = \cNP$, easy search, completeness, intermediate problems]
  \label{def:model-classes}
  \lean{Sagitta.Arrow.Model.PeqNP, Sagitta.Arrow.Model.SearchEasy, Sagitta.Arrow.Model.NPComplete, Sagitta.Arrow.Model.Intermediate}
  \leanok
  \uses{def:model}
  In a model $M$:
  $\cP = \cNP$ means that every $A \in \cNP$ is in $\cP$;
  \emph{search is easy} means that every $R \in \cFNP$ is in $\cFP$;
  $B$ is \emph{$\cNP$-complete} when $B \in \cNP$ and $A \red B$ for every $A \in \cNP$;
  $A$ is \emph{$\cNP$-intermediate} when $A \in \cNP$, $A \notin \cP$ and $A$ is not
  $\cNP$-complete.
  \own{}
\end{definition}

\section{Imported from cslib}
\label{sec:cslib}

The entries of this section are \textbf{not this project's work}. They are definitions and
theorems of \href{https://github.com/leanprover/cslib}{cslib} (commit \texttt{98f9a62}) and
\href{https://github.com/cslib-community/cslib-frontier}{cslib-frontier} (commit \texttt{99f0281}),
copied into this repository without modification, under their original paths and module names,
with their Apache-2.0 licenses and copyright headers (see
\href{https://github.com/julatec/sagitta/blob/main/VENDORED.md}{\texttt{VENDORED.md}}). Each entry
says where the copied file is. They are listed here so that the dependency graph shows on which
part of cslib each result of this project rests.

All problems are over instances $a$ of a type $\alpha$ with an injective binary encoding
$\enc : \alpha \hookrightarrow \{0,1\}^*$.

\subsection{The machines}

\begin{definition}[Multi-tape Turing machines]
  \label{def:cslib-tm}
  \lean{Turing.MultiTapeTM, Turing.Cfg, Turing.Action, Turing.MultiTapeTM.runFrom}
  \leanok
  A deterministic multi-tape machine with $k$ work tapes has a transition that, from a state, the
  symbol under the input head and the symbols under the work heads, gives an action: move the input
  head, write and move on each work tape, append at most one symbol to the output, and go to a new
  state or halt. A configuration records the state, the head positions, the work tapes and the
  output; $\mathrm{runFrom}(c, t)$ is the configuration after $t$ steps.
  \cslibfile{Cslib/Computability/Machines/Turing/MultiTape/Deterministic.lean}
  (configurations and actions in
  \cslibfile{Cslib/Computability/Machines/Turing/MultiTape/Configuration.lean}).
  \notours{cslib}
\end{definition}

\begin{definition}[Time and space of a run]
  \label{def:cslib-time-space}
  \lean{Turing.MultiTapeNTM.RunsInTime, Turing.MultiTapeNTM.RunsInSpace, Turing.MultiTapeTM.ComputesFunInTimeAndSpace}
  \leanok
  \uses{def:cslib-tm}
  A machine runs in time $t$ on an input when every run halts within $t$ steps, and in space $s$
  when it stays within $s$ cells; it computes $f$ in time $t$ and space $s$ (for encodings
  $\enc_{\mathrm{in}}, \enc_{\mathrm{out}}$) when on $\enc_{\mathrm{in}}(a)$ it halts with output
  $\enc_{\mathrm{out}}(f(a))$ within those bounds.
  \cslibfile{Cslib/Computability/Machines/Turing/MultiTape/Nondeterministic.lean},
  \cslibfile{Cslib/Computability/Machines/Turing/MultiTape/Space.lean} and
  \cslibfile{Cslib/Computability/Machines/Turing/MultiTape/Deterministic.lean}.
  \notours{cslib}
\end{definition}

\begin{definition}[Bounded configurations]
  \label{def:cslib-configbound}
  \lean{Turing.Cfg.core, Turing.Storage, Turing.storageBound}
  \leanok
  \uses{def:cslib-tm}
  The core of a configuration forgets the output; a configuration whose work tapes fit in windows of
  size $s$ is one of at most $\mathrm{storageBound}$ storages, a number exponential in $s$.
  \cslibfile{Cslib/Computability/Machines/Turing/MultiTape/ConfigBound.lean}.
  \notours{cslib}
\end{definition}

\begin{definition}[Plumbing: sequential composition, virtual input and output, rewinding]
  \label{def:cslib-plumbing}
  \lean{Turing.MultiTapeTM.seq, Turing.MultiTapeTM.inputFromTape, Turing.MultiTapeTM.outputToTape, Turing.MultiTapeTM.rewindWork}
  \leanok
  \uses{def:cslib-tm}
  $\mathrm{seq}$ runs one machine and then another on the same tapes; $\mathrm{inputFromTape}$ makes
  a machine read its input from a work tape; $\mathrm{outputToTape}$ writes its output on a fresh
  work tape; $\mathrm{rewindWork}$ brings a work head back to the start of its word.
  \cslibfile{Cslib/Computability/Machines/Turing/MultiTape/Plumbing/Sequential.lean},
  \cslibfile{Cslib/Computability/Machines/Turing/MultiTape/Plumbing/InputFromTape.lean},
  \cslibfile{Cslib/Computability/Machines/Turing/MultiTape/Plumbing/OutputToTape.lean},
  \cslibfile{Cslib/Computability/Machines/Turing/MultiTape/Plumbing/RewindWork.lean}.
  \notours{cslib}
\end{definition}

\subsection{The classes}

\begin{definition}[Polynomial-time computable function]
  \label{def:polytime}
  \lean{Turing.MultiTapeTM.ComputableInPolyTime}
  \leanok
  \uses{def:cslib-time-space}
  $f : \alpha \to \beta$ is computable in polynomial time when some machine with binary alphabet and
  finitely many states computes it within time and space bounded by polynomially bounded functions
  of the length of the encoded input.
  \cslibfile{CslibFrontier/Computability/Machines/Turing/MultiTape/PolyTime.lean}.
  \notours{cslib-frontier}
\end{definition}

\begin{definition}[The class $\cP$]
  \label{def:P}
  \lean{Cslib.Complexity.P}
  \leanok
  \uses{def:polytime}
  $A \in \cP$ when some Boolean function computable in polynomial time decides membership in $A$.
  \cslibfile{CslibFrontier/Computability/Complexity/Classes/P.lean}.
  \notours{cslib-frontier}
\end{definition}

\begin{definition}[The class $\cNP$, by verifiers]
  \label{def:NP}
  \lean{Cslib.Complexity.IsVerifier, Cslib.Complexity.NP}
  \leanok
  \uses{def:polytime}
  $V$ verifies $A$ with polynomial $p$ when
  $a \in A \iff \exists w,\ |w| = p(|\enc(a)|) \wedge V(a, w) = \mathrm{true}$; $A \in \cNP$ when it
  has a verifier computable in polynomial time.
  \cslibfile{CslibFrontier/Computability/Complexity/Classes/NP.lean}.
  \notours{cslib-frontier}
\end{definition}

\begin{definition}[Polynomial-time reductions and $\cNP$-completeness]
  \label{def:reduction}
  \lean{Cslib.Complexity.PolyTimeReducible, Cslib.Complexity.PolyTimeReducible.trans, Cslib.Complexity.NPComplete, Cslib.Complexity.P_subset_NP}
  \leanok
  \uses{def:P, def:NP}
  $A \red B$ when some $f$ computable in polynomial time satisfies $a \in A \iff f(a) \in B$;
  reductions compose. $B$ is $\cNP$-complete when $B \in \cNP$ and every $\cNP$ problem reduces to
  it. $\cP \subseteq \cNP$.
  \cslibfile{CslibFrontier/Computability/Complexity/Reduction.lean}.
  \notours{cslib-frontier}
\end{definition}

\begin{definition}[Search problems: $\cFNP$ and $\cFP$]
  \label{def:FNP}
  \lean{Cslib.Complexity.FNP, Cslib.Complexity.FPSearch}
  \leanok
  \uses{def:polytime}
  A search problem is a relation $R(a, w)$. $R \in \cFNP$ when its solutions are polynomially short
  and it is decided in polynomial time; $R \in \cFP$ when moreover a function computable in
  polynomial time returns a solution whenever there is one.
  \cslibfile{CslibFrontier/Computability/Complexity/Classes/TFNP.lean}.
  \notours{cslib-frontier}
\end{definition}

\begin{theorem}[Search reduces to decision]
  \label{thm:cslib-search-decision}
  \lean{Cslib.Complexity.FNP_subset_FPSearch_iff}
  \leanok
  \uses{def:FNP, def:P, def:NP}
  $\cFNP \subseteq \cFP$ for every encoding if and only if $\cNP \subseteq \cP$ for every encoding.
  \cslibfile{CslibFrontier/Computability/Complexity/Classes/TFNP/SearchToDecision.lean}.
  \notours{cslib-frontier}
\end{theorem}

\begin{proof}
  \leanok
  In cslib-frontier.
\end{proof}

\subsection{The problems}

\begin{definition}[$\SAT$ and $\TSAT$]
  \label{def:threesat}
  \lean{Cslib.Complexity.SAT, Cslib.Complexity.encCNF, Cslib.Complexity.ThreeSAT}
  \leanok
  $\SAT$ is the set of satisfiable CNF formulas over $\mathbb{N}$, with a binary encoding of the
  clause list; $\TSAT$ the satisfiable ones with at most three literals per clause.
  \cslibfile{CslibFrontier/Computability/Complexity/Problems/SAT.lean},
  \cslibfile{CslibFrontier/Computability/Complexity/Problems/ThreeSAT.lean}.
  \notours{cslib-frontier}
\end{definition}

\begin{theorem}[Cook--Levin]
  \label{thm:cook-levin}
  \lean{Cslib.Complexity.CookLevin.SAT_npComplete, Cslib.Complexity.ThreeSAT_npComplete, Cslib.Complexity.polyTimeReducible_SAT_ThreeSAT}
  \leanok
  \uses{def:threesat, def:reduction}
  $\SAT$ and $\TSAT$ are $\cNP$-complete, and $\SAT \red \TSAT$.
  \cslibfile{CslibFrontier/Computability/Complexity/CookLevin/Completeness.lean},
  \cslibfile{CslibFrontier/Computability/Complexity/Problems/ThreeSAT.lean}.
  \notours{cslib-frontier}
\end{theorem}

\begin{proof}
  \leanok
  In cslib-frontier.
\end{proof}

\begin{theorem}[$\MAXCUT$ is $\cNP$-complete]
  \label{thm:maxcut}
  \lean{Cslib.Complexity.MaxCut, Cslib.Complexity.MaxCut_npComplete}
  \leanok
  \uses{def:reduction}
  The weighted maximum-cut problem (is there a cut of weight at least $k$?) is $\cNP$-complete.
  \cslibfile{CslibFrontier/Computability/Complexity/Problems/MaxCut.lean}.
  \notours{cslib-frontier}
\end{theorem}

\begin{proof}
  \leanok
  In cslib-frontier.
\end{proof}

\begin{theorem}[Ladner]
  \label{thm:ladner}
  \lean{Cslib.Complexity.Ladner.exists_NP_not_P_not_npComplete}
  \leanok
  \uses{def:P, def:NP, def:reduction}
  If $\cP \neq \cNP$ for some encoding, some problem over binary words is in $\cNP$, not in $\cP$
  and not $\cNP$-complete.
  \cslibfile{CslibFrontier/Computability/Complexity/Classes/Ladner.lean}.
  \notours{cslib-frontier}
\end{theorem}

\begin{proof}
  \leanok
  In cslib-frontier.
\end{proof}

\section{Turing machines form a model}

\begin{definition}[The Turing-machine model]
  \label{def:tmModel}
  \lean{Sagitta.Arrow.TM.tmModel, Sagitta.Arrow.TM.Problem, Sagitta.Arrow.TM.SearchProblem}
  \leanok
  \uses{def:model, def:P, def:NP, def:reduction, def:FNP, thm:cslib-search-decision}
  Decision problems are triples $(\alpha, \enc, A)$, search problems triples $(\alpha, \enc, R)$,
  reductions are polynomial-time many-one reductions, and the classes are those of cslib-frontier.
  Every closure property of Definition~\ref{def:model} is a theorem of cslib-frontier, used with its
  proof: composition of reductions, closure of $\cP$ under reductions, the domain of an $\cFNP$
  relation in $\cNP$, every $\cNP$ problem the domain of its verifier relation, and search through
  prefixes.
\end{definition}

\begin{lemma}[Completeness in the two vocabularies]
  \label{lem:npComplete-iff}
  \lean{Sagitta.Arrow.TM.npComplete_iff}
  \leanok
  \uses{def:tmModel, def:model-classes}
  For a problem $(\beta, \enc, B)$, $\cNP$-completeness in the abstract model coincides with
  cslib's $\cNP$-completeness.
\end{lemma}

\begin{proof}
  \leanok
  Unfold both definitions; an $\cNP$ problem of the model is an $\cNP$ problem over some encoded type.
\end{proof}

\chapter{The laws}
\label{chap:laws}

The essay states three laws of computation, by analogy with thermodynamics. The zeroth and the
first are theorems; the second is a named hypothesis in two strengths.

\section{The zeroth law}

\begin{theorem}[Zeroth law]
  \label{thm:lawZero}
  \lean{Sagitta.Arrow.Model.lawZero}
  \leanok
  \uses{def:model}
  In every model, reductions are transitive: if $A \red B$ and $B \red C$ then $A \red C$.
  \own{} (the name and the reading as a law; transitivity itself is standard,
  \cite{AroraBarak}, Chapter~2).
\end{theorem}

\begin{proof}
  \leanok
  It is the closure property ``reductions compose'' of the model.
\end{proof}

\begin{corollary}[Zeroth law for Turing machines]
  \label{cor:lawZero-tm}
  \lean{Sagitta.Arrow.TM.lawZero, Cslib.Complexity.PolyTimeReducible.trans}
  \leanok
  \uses{thm:lawZero, def:tmModel, def:reduction}
  Polynomial-time many-one reductions between encoded problems compose.
\end{corollary}

\begin{proof}
  \leanok
  Theorem~\ref{thm:lawZero} applied to the Turing-machine model; the field is cslib-frontier's
  \texttt{PolyTimeReducible.trans}.
\end{proof}

\begin{lemma}[Completeness travels along reductions]
  \label{lem:npcomplete-reduces}
  \lean{Sagitta.Arrow.Model.NPComplete.of_reduces, Sagitta.Arrow.Model.NPComplete.InP_iff}
  \leanok
  \uses{thm:lawZero, def:model-classes}
  If $B$ is $\cNP$-complete, $C \in \cNP$ and $B \red C$, then $C$ is $\cNP$-complete. An
  $\cNP$-complete $B$ is in $\cP$ if and only if $\cP = \cNP$.
\end{lemma}

\begin{proof}
  \leanok
  Compose the reductions; for the second part use the closure of $\cP$ under reductions.
\end{proof}

\section{The first law}

A \emph{query tree} over a finite set of options $\alpha$ asks, at each node, whether an option is
good, and answers yes or no at the leaves; its cost on an oracle $f : \alpha \to \{0,1\}$ is the
number of questions on the path that $f$ selects.

\begin{definition}[Query trees that decide existence]
  \label{def:querytree}
  \lean{Sagitta.Arrow.QueryTree, Sagitta.Arrow.QueryTree.cost, Sagitta.Arrow.QueryTree.DecidesExists}
  \leanok
  A query tree $t$ \emph{decides existence} when, for every oracle $f$, it answers yes exactly when
  $f(i) = 1$ for some option $i$.
\end{definition}

\begin{lemma}[The adversary]
  \label{lem:adversary}
  \lean{Sagitta.Arrow.QueryTree.adversary, Sagitta.Arrow.QueryTree.cost_adversary}
  \leanok
  \uses{def:querytree}
  A query tree that decides existence asks, on the all-zero oracle, about every option; so its cost
  there is at least $|\alpha|$.
  \source{the adversary argument for unstructured search; see \cite{AroraBarak}, Chapter~12
  (decision trees).}
\end{lemma}

\begin{proof}
  \leanok
  If an option $i$ is never asked on the all-zero oracle, the oracle that is $1$ only at $i$ follows
  the same path and gets the same answer, which is wrong for one of the two.
\end{proof}

\begin{theorem}[First law: work is conserved in a black box]
  \label{thm:conservation}
  \lean{Sagitta.Arrow.DependsOnlyOnAsked, Sagitta.Arrow.conservation}
  \leanok
  \uses{lem:adversary}
  Let a filling phase build a table $\mathrm{fill}(f)$ asking the options $\mathrm{asked}(f)$, and let
  the table depend only on the answers to those options (the system is closed). If a query phase
  $Q(\mathrm{fill}(f))$, reading the table for free, decides whether some option is good for every
  oracle $f$, then on the all-zero oracle
  \[
    |\alpha| \le |\mathrm{asked}(0)| + \mathrm{cost}(Q(\mathrm{fill}(0)), 0).
  \]
  What is hard is neither created nor destroyed: it is only moved between the phases.
  \own{} (the statement as a conservation law; the proof is the adversary argument).
\end{theorem}

\begin{proof}
  \leanok
  \uses{lem:adversary}
  Every option is asked by one of the two phases on the all-zero oracle; otherwise flip it and both
  phases behave the same.
\end{proof}

\begin{proposition}[Control: information from outside breaks the first law]
  \label{prop:free-table}
  \lean{Sagitta.Arrow.free_table, Sagitta.Arrow.free_table_not_closed, Sagitta.Arrow.free_table_breaks}
  \leanok
  \uses{thm:conservation}
  If the table is the oracle itself, the query phase decides at cost $0$ and the filling phase asks
  nothing; the table then does not depend only on what was asked, and the inequality of
  Theorem~\ref{thm:conservation} fails.
  \own{}
\end{proposition}

\begin{proof}
  \leanok
  Direct computation.
\end{proof}

\section{The second law}

\begin{definition}[The arrow, law 2.1]
  \label{def:arrow}
  \lean{Sagitta.Arrow.Model.Arrow}
  \leanok
  \uses{def:model-classes}
  The \emph{arrow} of a model is $\neg(\cP = \cNP)$, the worst-case version. It is a
  \textbf{hypothesis}: a definition of type \texttt{Prop}, taken as an explicit argument by every
  theorem that uses it, never an axiom.
  \own{} (the reading as a second law; the statement is $\cP \neq \cNP$).
\end{definition}

\begin{theorem}[Kelvin's form]
  \label{thm:kelvin}
  \lean{Sagitta.Arrow.Model.kelvin}
  \leanok
  \uses{def:arrow, thm:searchEasy-iff}
  The arrow holds if and only if search is not easy: there is no efficient process whose only
  effect is to convert every ``yes'' into its witness.
  \own{}
\end{theorem}

\begin{proof}
  \leanok
  \uses{thm:searchEasy-iff}
  The localization, part (3): search is easy exactly when $\cP = \cNP$.
\end{proof}

\begin{definition}[Two models]
  \label{def:two-models}
  \lean{Sagitta.Arrow.collapsedModel, Sagitta.Arrow.separatedModel}
  \leanok
  \uses{def:model}
  The \emph{collapsed} model has one problem, in $\cP$, and every search problem easy; the
  \emph{separated} model satisfies the closure properties with $\cP \neq \cNP$.
  \own{}
\end{definition}

\begin{theorem}[The arrow is independent of the closure properties]
  \label{thm:arrow-independent}
  \lean{Sagitta.Arrow.arrow_independent}
  \leanok
  \uses{def:arrow, def:two-models}
  The arrow fails in the collapsed model and holds in the separated model.
  \own{}
\end{theorem}

\begin{proof}
  \leanok
  Compute in both models.
\end{proof}

\begin{definition}[The arrow for Turing machines]
  \label{def:arrow-tm}
  \lean{Sagitta.Arrow.TM.Arrow}
  \leanok
  \uses{def:arrow, def:tmModel}
  The arrow of the Turing-machine model: for some encoded type, $\cNP \not\subseteq \cP$. A
  hypothesis, never an axiom.
\end{definition}

\begin{theorem}[The arrow, unfolded]
  \label{thm:arrow-iff-tm}
  \lean{Sagitta.Arrow.TM.arrow_iff, Sagitta.Arrow.TM.arrow_iff_exists, Sagitta.Arrow.TM.arrow_iff_SAT, Sagitta.Arrow.TM.kelvin}
  \leanok
  \uses{def:arrow-tm, thm:kelvin, thm:cook-levin, thm:oneBit}
  For Turing machines the arrow is equivalent to each of: not every $\cNP$ class is contained in
  $\cP$; $\cP \neq \cNP$ for some encoding; $\SAT \notin \cP$; and (Kelvin's form) not every $\cFNP$
  search problem is in $\cFP$.
\end{theorem}

\begin{proof}
  \leanok
  \uses{thm:cook-levin, thm:oneBit}
  The third equivalence is the Cook--Levin theorem; the fourth is cslib-frontier's
  search-to-decision theorem.
\end{proof}

\begin{definition}[Law 2.2]
  \label{def:law22}
  \lean{Sagitta.Arrow.TM.InTIME, Sagitta.Arrow.TM.subexpBound, Sagitta.Arrow.TM.Law22}
  \leanok
  \uses{def:tmModel, def:threesat}
  \emph{Law 2.2}: for every $\varepsilon > 0$, $\TSAT \notin \TIME(2^{n^{1-\varepsilon}})$, where
  $\TIME(T)$ allows $C \cdot T(n) + C$ steps for a constant $C$ and $n$ is the length of the encoded
  formula. A hypothesis, never an axiom.
  \source{the weak exponential-time hypothesis, shown equivalent to Stearns and Hunt's satisfiability
  hypothesis SH in \cite{ETHdead} (as cited in the Lean source, Theorem~2.4).}
  \fidelity{$n$ is measured in bits of the CNF encoding (the size form of SH), not in variables. It is
  equivalent to the weak ETH ($n$ variables) only by Theorem~2.4 of \cite{ETHdead}, which is not
  formalized here.}
\end{definition}

\begin{theorem}[Law 2.2 implies law 2.1]
  \label{thm:law22-arrow}
  \lean{Sagitta.Arrow.TM.Law22.arrow, Sagitta.Arrow.TM.poly_le_subexp}
  \leanok
  \uses{def:law22, def:arrow-tm, def:threesat}
  Law 2.2 implies the arrow.
\end{theorem}

\begin{proof}
  \leanok
  If $\cP = \cNP$ then $\TSAT$ is decided in polynomial time, hence in time
  $C \cdot 2^{n^{1/2}} + C$, which law 2.2 forbids with $\varepsilon = 1/2$.
\end{proof}

\chapter{The localization}
\label{chap:localization}

Where does the asymmetry between the two directions of a search live? The localization answers
in three parts: \textbf{not} in the dynamics (part 1), \textbf{not} in the state while it is kept
whole (part 2), but \textbf{only} in the erasure of the boundary down to one bit (part 3), where it
is exactly $\cP$ versus $\cNP$. The whole chapter is this project's own; the ingredients that come
from the literature are named in each entry.

\section{Part 1: the dynamics does not carry the arrow}

\subsection{Reversible dynamics}

\begin{definition}[Reversible dynamics]
  \label{def:reversible-dynamics}
  \lean{Sagitta.Arrow.ReversibleDynamics}
  \leanok
  A \emph{reversible dynamics} on a type $\sigma$ is a step $s : \sigma \to \sigma$ with a left
  inverse $b$ (so $b(s(x)) = x$), both of the same cost per step.
  \own{}
\end{definition}

\begin{theorem}[Going back costs the same as going forward]
  \label{thm:back-iterate}
  \lean{Sagitta.Arrow.ReversibleDynamics.back_iterate_step_iterate, Sagitta.Arrow.ReversibleDynamics.costBack_eq_costForward, Sagitta.Arrow.ReversibleDynamics.step_injective}
  \leanok
  \uses{def:reversible-dynamics}
  For every $n$ and $x$, $b^n(s^n(x)) = x$, and the cost of $n$ steps back equals the cost of $n$
  steps forward. The step is injective.
  \own{}
\end{theorem}

\begin{proof}
  \leanok
  Induction on $n$.
\end{proof}

\begin{theorem}[Any dynamics becomes reversible with a history]
  \label{thm:history-dynamics}
  \lean{Sagitta.Arrow.historyDynamics, Sagitta.Arrow.history_simulates, Sagitta.Arrow.history_recovers, Sagitta.Arrow.injective_iff_exists_back}
  \leanok
  \uses{def:reversible-dynamics, thm:back-iterate}
  For any $f : \sigma \to \sigma$, the dynamics on $\sigma \times \sigma^*$ that applies $f$ and pushes
  the old state on a history list is reversible, simulates $f$, and recovers the initial state after
  $n$ steps back. A map has a left inverse if and only if it is injective.
  \source{this is the abstract shadow of Bennett's construction \cite{Bennett73}; the
  Turing-machine version is Chapter~\ref{chap:bennett}.}
\end{theorem}

\begin{proof}
  \leanok
  Pop the history.
\end{proof}

\subsection{Time reversal: information is conserved (discrete Noether)}

\begin{definition}[Koopman operator]
  \label{def:koopman}
  \lean{Sagitta.Arrow.Dynamics.koopman, Sagitta.Arrow.Dynamics.mult, Sagitta.Arrow.Dynamics.phase, Sagitta.Arrow.Dynamics.norm2}
  \leanok
  For $f : \sigma \to \sigma$ on a finite set, the Koopman operator is the linear map
  $K_f \varphi = \varphi \circ f$ on $\mathbb{C}^\sigma$; $M_\varphi$ is multiplication by $\varphi$,
  $e^{i\theta Q}$ the phase of a real quantity $Q$, and $\|\varphi\|^2 = \sum_x |\varphi(x)|^2$.
  \source{\cite{Koopman31}.}
\end{definition}

\begin{theorem}[Discrete Noether]
  \label{thm:noether}
  \lean{Sagitta.Arrow.Dynamics.conserved_iff_const_on_orbits, Sagitta.Arrow.Dynamics.koopman_noether}
  \leanok
  \uses{def:koopman}
  For a permutation $f$, a quantity $Q$ is conserved ($Q \circ f = Q$) if and only if it is constant
  on the orbits of the group generated by $f$. For any $f$ on a finite set, $Q \circ f = Q$ if and
  only if $K_f$ commutes with $M_{e^{i\theta Q}}$ for every $\theta$.
  \own{} (a discrete analogue of Noether's theorem in Koopman's picture).
\end{theorem}

\begin{proof}
  \leanok
  Orbits: induction on the powers of $f$. Koopman: evaluate the commutator on indicator functions and
  separate the values of $Q$ with a suitable $\theta$.
\end{proof}

\begin{theorem}[Unitarity is bijectivity]
  \label{thm:koopman-isometry}
  \lean{Sagitta.Arrow.Dynamics.koopman_isometry_iff, Sagitta.Arrow.Dynamics.erasure_not_isometry}
  \leanok
  \uses{def:koopman}
  $K_f$ preserves $\|\cdot\|^2$ if and only if $f$ is bijective. Erasure (the constant map on
  $\{0,1\}$) is not an isometry.
  \own{}
\end{theorem}

\begin{proof}
  \leanok
  Evaluate the norm on indicator functions of single states and of sets.
\end{proof}

\begin{theorem}[Reversible dynamics neither creates nor destroys information]
  \label{thm:card-image}
  \lean{Sagitta.Arrow.ReversibleDynamics.card_image_iterate, Sagitta.Arrow.ReversibleDynamics.koopman_isometry, Sagitta.Arrow.Dynamics.entropy_pushforward}
  \leanok
  \uses{def:reversible-dynamics, thm:koopman-isometry}
  A reversible dynamics keeps the size of every finite set of states after $n$ steps, its Koopman
  operator is an isometry on a finite type, and a permutation preserves the Shannon entropy of a
  distribution.
  \own{}
\end{theorem}

\begin{proof}
  \leanok
  The step is injective (Theorem~\ref{thm:back-iterate}).
\end{proof}

\subsection{Time translation: the system does not know the time}

\begin{definition}[Homogeneous automaton]
  \label{def:automaton}
  \lean{Sagitta.Arrow.Dynamics.Automaton, Sagitta.Arrow.Dynamics.Automaton.run, Sagitta.Arrow.Dynamics.Automaton.step}
  \leanok
  An automaton that consults a verifier has states $\sigma$, an initial state, a query $q(x)$ and an
  output $o(x) \in Y \cup \{\bot\}$ for each state, and a transition $\delta(x, r)$ on the answer
  $r$. It is \emph{homogeneous in time}: there is no step counter, and what it knows about the time
  is what it keeps in $\sigma$.
  \own{}
\end{definition}

\begin{theorem}[The Noether charge of time translation]
  \label{thm:noether-charge}
  \lean{Sagitta.Arrow.Dynamics.noether_charge, Sagitta.Arrow.Dynamics.does_not_know_the_time}
  \leanok
  \uses{def:automaton}
  Two runs against the same verifier that reach the same state, one at time $i$ and the other at
  time $j$, without having stopped, continue identically. In particular a run that comes back to a
  state continues from $i$ and from $j$ in the same way: it does not know the time.
  \own{}
\end{theorem}

\begin{proof}
  \leanok
  Split the run at $i$ and at $j$.
\end{proof}

\begin{definition}[Clocked automaton]
  \label{def:clocked}
  \lean{Sagitta.Arrow.Dynamics.ClockedAutomaton, Sagitta.Arrow.Dynamics.ClockedAutomaton.run, Sagitta.Arrow.Dynamics.ClockedAutomaton.homogenize}
  \leanok
  \uses{def:automaton}
  A clocked automaton also reads a time register $\kappa$ that advances on its own. Its
  \emph{homogenization} is the homogeneous automaton on $\sigma \times \kappa$.
  \own{}
\end{definition}

\begin{theorem}[A clock is paid in memory]
  \label{thm:homogenize}
  \lean{Sagitta.Arrow.Dynamics.ClockedAutomaton.homogenize_run, Sagitta.Arrow.Dynamics.card_clock_bits}
  \leanok
  \uses{def:clocked}
  The homogenization has the same runs as the clocked automaton, and a clock of $c$ bits multiplies
  the number of states by $2^c$.
  \own{}
\end{theorem}

\begin{proof}
  \leanok
  Induction on the run.
\end{proof}

\begin{theorem}[The cube law]
  \label{thm:cube-law}
  \lean{Sagitta.Arrow.Dynamics.cube_law, Sagitta.Arrow.Dynamics.cube_law_general, Sagitta.Arrow.Dynamics.cube_law_clock}
  \leanok
  \uses{def:automaton, thm:homogenize}
  If automata indexed by at most $2^L$ programs and $2^s$ advice strings, with at most $2^S$ states
  each, find any of $N$ targets with $t$ yes/no queries, then
  $N \le 2^L \cdot 2^s \cdot \min(t+1, 2^{S-1})$. With answers of $b$ bits the bound is
  $2^L \cdot 2^s \cdot \min(2^{bt}, 2^S)$; with a clock of $c$ bits, $S$ becomes $S + c$.
  \own{}
\end{theorem}

\begin{proof}
  \leanok
  Count the targets found along the all-``no'' trajectory, bounded both by the time and by the
  number of states, and take a union over programs and advice.
\end{proof}

\begin{theorem}[Halt or cycle; the CHT with its costs]
  \label{thm:cht}
  \lean{Sagitta.Arrow.Dynamics.halt_or_cycle, Sagitta.Arrow.Dynamics.halts_before_two_pow, Sagitta.Arrow.Dynamics.cht_time, Sagitta.Arrow.Dynamics.clock_in_memory}
  \leanok
  \uses{thm:noether-charge, thm:homogenize}
  A deterministic system on a finite set either halts within $|\sigma|$ steps or never halts and is
  eventually periodic. Hence a homogeneous automaton with $|\sigma| \le 2^S$ that answers, answers
  in fewer than $2^S$ queries; with a clock of $c$ bits, in fewer than $2^{S+c}$.
  \own{} (``CHT'' is this project's name for the bound; the pigeonhole argument is classical).
\end{theorem}

\begin{proof}
  \leanok
  Pigeonhole on the trace, and the Noether charge after a repeated state.
\end{proof}

\begin{theorem}[The CHT bound is exponential, and attained]
  \label{thm:cht-exponential}
  \lean{Sagitta.Arrow.Dynamics.counter, Sagitta.Arrow.Dynamics.cht_exponential_not_polynomial}
  \leanok
  \uses{thm:cht}
  For every $n \ge 1$ there is a homogeneous automaton with $n+1$ bits of memory (a counter) that
  always answers within the budget $2^{n+1}$ but does not answer in $2^n - 1$ queries.
  \own{}
\end{theorem}

\begin{proof}
  \leanok
  The counter walks through $2^n$ states before it may answer.
\end{proof}

\subsection{Part 1 for Turing machines}

\begin{definition}[A machine with bounded space is a homogeneous automaton]
  \label{def:toAutomaton}
  \lean{Sagitta.Arrow.TM.toAutomaton, Sagitta.Arrow.TM.spaceBits}
  \leanok
  \uses{def:automaton, def:cslib-tm, def:cslib-time-space, def:cslib-configbound}
  A cslib multi-tape machine with $k$ work tapes, run on an input of length $n$ in space $S$, is a
  homogeneous automaton whose states are its configurations without the output (the state, the
  input head and the windows of the work tapes), at most $2^{\mathrm{spaceBits}}$ of them.
  \own{}
\end{definition}

\begin{theorem}[cslib's machines do not know the time]
  \label{thm:tm-dynamics}
  \lean{Sagitta.Arrow.TM.run_toAutomaton, Sagitta.Arrow.TM.card_states_bits, Sagitta.Arrow.TM.does_not_know_the_time_tm, Sagitta.Arrow.TM.cube_law_tm, Sagitta.Arrow.TM.halt_or_cycle_tm, Sagitta.Arrow.TM.halts_soon_tm}
  \leanok
  \uses{def:toAutomaton, thm:noether-charge, thm:cube-law, thm:cht}
  A machine that runs in space $S$ on an input does not know the time (if its configuration repeats
  before it halts, the run repeats), obeys the cube law, halts or cycles, and if it halts, halts
  within $2^{\mathrm{spaceBits}}$ steps.
  \own{}
\end{theorem}

\begin{proof}
  \leanok
  The automaton of Definition~\ref{def:toAutomaton} runs in lockstep with the machine.
\end{proof}

\begin{theorem}[CHT tower, piece 1: the $\cPSPACE$ bound]
  \label{thm:cht-pspace}
  \lean{Sagitta.Arrow.TM.cht_pspace}
  \leanok
  \uses{thm:tm-dynamics}
  A machine that runs in polynomial space and halts on an input halts within $2^{b(n)}$ steps for a
  polynomially bounded $b$.
  \source{the classical bound $\mathsf{SPACE}(s) \subseteq \TIME(2^{O(s)})$; \cite{AroraBarak},
  Chapter~4.}
\end{theorem}

\begin{proof}
  \leanok
  Theorem~\ref{thm:tm-dynamics}, with $\mathrm{spaceBits}$ polynomial in $n$.
\end{proof}

\begin{theorem}[Localization, part 1, complete]
  \label{thm:localization-one}
  \lean{Sagitta.Arrow.Dynamics.localization_one, Sagitta.Arrow.TM.localization_one_tm}
  \leanok
  \uses{thm:noether, thm:koopman-isometry, thm:card-image, thm:noether-charge, thm:homogenize, thm:cht, thm:arrow-independent, thm:hd-iff-arrow, thm:tm-dynamics, thm:cht-pspace, thm:cht-exponential, thm:hd-iff-arrow-tm}
  Neither symmetry of time decides the arrow. Time reversal conserves information (discrete Noether,
  unitarity iff bijectivity, sizes of sets kept); time translation makes a homogeneous system blind
  to the clock and charges the clock to memory, and the CHT bounds by $2^S$ the time of a system of
  $S$ bits. These facts hold in every model with no hypothesis, while the arrow holds in one model
  and fails in another; what decides it is the decision hypothesis, which is about the bit, not the
  dynamics. The same holds for cslib's machines.
  \own{}
  \fidelity{a conjunction of the proved facts. ``Neither symmetry decides the arrow'' is the essay's
  reading, not a formal statement; \texttt{arrow\_independent} uses models without dynamics.}
\end{theorem}

\begin{proof}
  \leanok
  The conjunction of the cited results.
\end{proof}

\section{Part 2: the full boundary, and the interpolation}

\begin{definition}[Completing a prefix]
  \label{def:complete}
  \lean{Sagitta.Arrow.complete}
  \leanok
  Given a verifier $V$ on bit strings and a prefix $u$, $\mathrm{complete}(V, u, k)$ searches the
  $2^k$ completions of $u$ by depth-first search and returns the first $w$ with $V(w) = 1$, together
  with the number of calls to $V$.
  \own{}
\end{definition}

\begin{theorem}[The full boundary]
  \label{thm:full-boundary}
  \lean{Sagitta.Arrow.full_boundary}
  \leanok
  \uses{def:complete}
  When the whole witness $w$ is kept ($V(w) = 1$), going back is one verification:
  $\mathrm{complete}(V, w, 0) = (w, 1)$.
  \own{}
\end{theorem}

\begin{proof}
  \leanok
  By definition.
\end{proof}

\begin{theorem}[The interpolation]
  \label{thm:interpolation}
  \lean{Sagitta.Arrow.interpolation, Sagitta.Arrow.complete_calls_le, Sagitta.Arrow.complete_sound, Sagitta.Arrow.complete_isSome}
  \leanok
  \uses{def:complete}
  With $m$ of the $n$ bits of a witness kept (a prefix $u$ of length $m$ that has a good completion
  of length $n$), going back finds a witness of length $n$ extending $u$ with at most $2^{n-m}$ calls
  to the verifier, so at cost at most $2^{n-m} c$ when one call costs $c$.
  \own{}
\end{theorem}

\begin{proof}
  \leanok
  Induction on the number of missing bits: each level doubles the calls at most.
\end{proof}

\begin{theorem}[The interpolation is tight in a black box]
  \label{thm:boundary-lower-bound}
  \lean{Sagitta.Arrow.boundary_lower_bound}
  \leanok
  \uses{lem:adversary}
  Any query tree over the $2^k$ completions that decides whether one of them is good asks, on the
  all-zero oracle, about all $2^k$.
  \own{}
\end{theorem}

\begin{proof}
  \leanok
  \uses{lem:adversary}
  The adversary lemma with $|\alpha| = 2^k$.
\end{proof}

\section{Part 3: one bit}

\begin{theorem}[With one bit, going back is cheap exactly when $\cP = \cNP$]
  \label{thm:searchEasy-iff}
  \lean{Sagitta.Arrow.Model.searchEasy_iff_peqNP}
  \leanok
  \uses{def:model-classes}
  In every model, search is easy if and only if $\cP = \cNP$.
  \source{search-to-decision \cite{Goldreich}, Chapter~2; here derived from the closure properties
  alone.}
\end{theorem}

\begin{proof}
  \leanok
  Decision reduces to search through the domain; search reduces to decision through the prefix
  problem, one call per bit.
\end{proof}

\begin{theorem}[One bit, for Turing machines]
  \label{thm:oneBit}
  \lean{Sagitta.Arrow.TM.oneBit, Sagitta.Arrow.TM.searchEasy_iff}
  \leanok
  \uses{thm:cslib-search-decision, def:tmModel}
  $\cFNP \subseteq \cFP$ for every encoding if and only if $\cNP \subseteq \cP$ for every encoding.
\end{theorem}

\begin{proof}
  \leanok
  \uses{thm:cslib-search-decision}
  It is cslib-frontier's theorem (Theorem~\ref{thm:cslib-search-decision}).
\end{proof}

\section{The localization theorem}

\begin{theorem}[Localization]
  \label{thm:localization}
  \lean{Sagitta.Arrow.localization, Sagitta.Arrow.TM.localizationTM}
  \leanok
  \uses{thm:back-iterate, thm:full-boundary, thm:interpolation, thm:searchEasy-iff, thm:oneBit}
  In every model, and for Turing machines:
  \begin{enumerate}
    \item a reversible dynamics, run back from its final state, recovers the initial state at the
      same cost as running it forward;
    \item with the full boundary going back is one verification, and with $m$ of the $n$ bits it is at
      most $2^{n-m}$ verifications;
    \item with one bit (only the answer ``yes'') going back is cheap for every search problem if and
      only if $\cP = \cNP$.
  \end{enumerate}
  \own{}
  \fidelity{a conjunction of the three parts. The word ``only'' is the essay's reading, not a formal
  statement. In the abstract core the equality of the cost back and forward is definitional; for
  Turing machines there is no general costed inverse.}
\end{theorem}

\begin{proof}
  \leanok
  The conjunction of the three parts.
\end{proof}

\chapter{The decision hypothesis}
\label{chap:decision}

The central hypothesis is not stated as $\cP \neq \cNP$ but at the place where the localization
says the asymmetry lives: reducing the final state to one bit is the boundary condition that breaks
the symmetry of time, and that breaking has a cost. Everything in this chapter is this project's
own.

\begin{definition}[The decision hypothesis, HD]
  \label{def:hd}
  \lean{Sagitta.Arrow.Model.DecisionHypothesis}
  \leanok
  \uses{def:model}
  In a model $M$, \emph{HD} says that some search problem $R \in \cFNP$ is not in $\cFP$: from the
  one bit ``the instance is in $\dom R$'' no efficient procedure reconstructs a witness. A
  \textbf{hypothesis}: a definition of type \texttt{Prop}, never an axiom.
  \own{}
\end{definition}

\begin{theorem}[HD is equivalent to the arrow]
  \label{thm:hd-iff-arrow}
  \lean{Sagitta.Arrow.Model.decisionHypothesis_iff_not_searchEasy, Sagitta.Arrow.Model.decisionHypothesis_iff_arrow, Sagitta.Arrow.Model.DecisionHypothesis.arrow}
  \leanok
  \uses{def:hd, def:arrow, thm:kelvin, thm:searchEasy-iff}
  In every model, HD holds if and only if search is not easy, if and only if the arrow holds. HD is
  the conceptual formulation of $\cP \neq \cNP$ (worst case), not a stronger hypothesis.
  \own{}
\end{theorem}

\begin{proof}
  \leanok
  \uses{thm:kelvin, thm:searchEasy-iff}
  HD is the negation of easy search, and Kelvin's form (through the localization, part 3) identifies
  that with the arrow.
\end{proof}

\begin{theorem}[Where HD lives, on the boundary]
  \label{thm:hd-boundary}
  \lean{Sagitta.Arrow.decisionHypothesis_on_boundary}
  \leanok
  \uses{thm:full-boundary, thm:boundary-lower-bound, thm:hd-iff-arrow}
  With the full boundary, going back is one verification; with nothing but the bit, a black-box
  search over the $2^k$ completions asks about all of them in the worst case; and whether a
  non-black-box procedure does better from that bit, for every verifier, is what HD denies, and HD
  is the arrow.
  \own{}
\end{theorem}

\begin{proof}
  \leanok
  The conjunction of the cited results.
\end{proof}

\begin{theorem}[HD is independent of the closure properties; the consequences need it]
  \label{thm:hd-independent}
  \lean{Sagitta.Arrow.decisionHypothesis_independent, Sagitta.Arrow.consequences_need_decisionHypothesis}
  \leanok
  \uses{def:hd, def:two-models, thm:hd-iff-arrow}
  HD fails in the collapsed model and holds in the separated one. In the collapsed model, where HD
  fails, the only problem is $\cNP$-complete and in $\cP$.
  \own{}
\end{theorem}

\begin{proof}
  \leanok
  Compute in both models.
\end{proof}

\begin{definition}[HD for Turing machines]
  \label{def:hd-tm}
  \lean{Sagitta.Arrow.TM.DecisionHypothesis}
  \leanok
  \uses{def:polytime, def:FNP}
  For some encoded type, some polynomial $p$ and some verifier $V$ computable in polynomial time, no
  function computable in polynomial time returns, from the instance alone, a certificate of length
  $p(n)$ accepted by $V$ whenever there is one. A hypothesis, never an axiom.
  \own{}
\end{definition}

\begin{theorem}[HD is equivalent to the arrow, for Turing machines]
  \label{thm:hd-iff-arrow-tm}
  \lean{Sagitta.Arrow.TM.decisionHypothesis_iff_arrow, Sagitta.Arrow.TM.DecisionHypothesis.arrow, Sagitta.Arrow.TM.decisionHypothesis_iff_model}
  \leanok
  \uses{def:hd-tm, def:arrow-tm, thm:oneBit, thm:hd-iff-arrow}
  For Turing machines, HD holds if and only if the arrow holds, and HD coincides with the HD of the
  Turing-machine model.
  \own{}
\end{theorem}

\begin{proof}
  \leanok
  \uses{thm:oneBit}
  The localization, part 3, for Turing machines.
\end{proof}

\begin{theorem}[Irreversibility without loss of information]
  \label{thm:irreversibility-without-loss}
  \lean{Sagitta.Arrow.Dynamics.irreversibility_without_loss, Sagitta.Arrow.TM.irreversibility_without_loss_tm}
  \leanok
  \uses{def:hd, def:hd-tm, thm:card-image, thm:hd-iff-arrow, thm:hd-iff-arrow-tm, thm:sat-notmem-p}
  Under HD, both hold at once: a reversible dynamics loses no information (it keeps the size of
  every set of states, and its Koopman operator is unitary --- a theorem), and yet from the one bit
  ``a witness exists'' the witness is not recovered cheaply for every search problem (HD, a
  hypothesis). For Turing machines, moreover $\SAT \notin \cP$. The irreversibility lives in the
  bit, not in the dynamics.
  \own{}
\end{theorem}

\begin{proof}
  \leanok
  The conjunction of Theorem~\ref{thm:card-image} with HD.
\end{proof}

\chapter{Bennett's theorem for cslib's machines}
\label{chap:bennett}

Bennett showed in 1973 that every computation can be simulated by a logically reversible one
\cite{Bennett73}: compute keeping a history of what each step forgets, copy the result, and
uncompute. Here the statement is proved for the multi-tape machines of cslib, one file of
\texttt{Sagitta/Arrow/TM/Bennett} per step of the construction. With it, the reversible
dynamics of Chapter~\ref{chap:localization} is not an assumption about machines: every
polynomial-time verifier runs on a reversible machine.

\section{The statement}

\begin{definition}[Reversible machine]
  \label{def:isReversible}
  \lean{Sagitta.Arrow.TM.IsReversible}
  \leanok
  \uses{def:cslib-tm}
  A machine is \emph{reversible} when its step is injective on running configurations: every
  running configuration has at most one running predecessor. The condition is on \emph{every}
  configuration, not only the reachable ones, and the configuration includes the output.
  \source{\cite{Bennett73}, logical reversibility.}
  \fidelity{Bennett's machines are quadruple machines on one or three tapes; here reversibility is
  injectivity of the step of a cslib multi-tape machine, with its read-only input and its
  append-only output.}
\end{definition}

\begin{lemma}[A reversible machine can be traced back]
  \label{lem:runFrom-injective}
  \lean{Sagitta.Arrow.TM.IsReversible.runFrom_injective}
  \leanok
  \uses{def:isReversible}
  If two runs of $n$ steps of a reversible machine do not halt before the end and reach the same
  configuration, they started from the same configuration.
\end{lemma}

\begin{proof}
  \leanok
  Induction on $n$, using injectivity of the step.
\end{proof}

\begin{definition}[Bennett's statement]
  \label{def:bennett}
  \lean{Sagitta.Arrow.TM.Bennett}
  \leanok
  \uses{def:isReversible, def:polytime}
  \emph{Bennett}: every function computable in polynomial time is computed in polynomial time and
  polynomial space by a reversible machine with finitely many states.
  \source{\cite{Bennett73}.}
  \fidelity{Bennett's theorem is stated for any computation, with time and space overheads; here it
  is restricted to polynomial-time functions, with polynomial bounds. The statement does not yet say
  that the auxiliary tapes end clean, which is the point of Bennett (1973): a history simulation
  that never uncomputes would satisfy it. Clean tapes are being added to the statement; see the
  Fidelity notes of the README.}
\end{definition}

\begin{lemma}[A reversible machine stays on a blank of its input]
  \label{lem:staysOnBlank}
  \lean{Sagitta.Arrow.TM.Bennett.StaysOnBlank, Sagitta.Arrow.TM.IsReversible.staysOnBlank}
  \leanok
  \uses{def:isReversible}
  A machine \emph{stays on a blank} when it never moves its input head while it reads a blank (the
  two boundary cells). Every reversible machine does: on the empty input, two configurations that
  differ only in the input head (on cell $0$ or $1$) take the same action, and any move sends both
  to the same cell.
  \own{} (a constraint that cslib's input-tape conventions impose on reversibility).
\end{lemma}

\begin{proof}
  \leanok
  Exhibit the two configurations and use injectivity.
\end{proof}

\section{Step 2: the machine with a history}

\begin{lemma}[Local reversibility of an action]
  \label{lem:apply-injective}
  \lean{Sagitta.Arrow.TM.Bennett.apply_injective}
  \leanok
  \uses{def:cslib-tm}
  The same action, applied to two configurations in the same state that read the same symbols where
  it writes, and whose input head does not move from a blank, gives the same result only if the two
  configurations were equal.
\end{lemma}

\begin{proof}
  \leanok
  Compare tape by tape.
\end{proof}

\begin{definition}[The history machine]
  \label{def:histTM}
  \lean{Sagitta.Arrow.TM.Bennett.Record, Sagitta.Arrow.TM.Bennett.histTM, Sagitta.Arrow.TM.Bennett.proj, Sagitta.Arrow.TM.Bennett.oneHot}
  \leanok
  \uses{def:cslib-tm}
  A \emph{record} is what a step forgets: the state and the symbols read. The history machine of
  $M$ has the $k$ tapes of $M$ and $h$ history tapes; on a clean history cell it takes the action of
  $M$ and writes the code of the record, one bit per history tape, moving the history heads right.
  On a configuration with garbage in the history (which no run reaches) it goes to a sink that
  appends to the output forever. The projection forgets the history tapes.
  \source{\cite{Bennett73}, the history tape.}
\end{definition}

\begin{theorem}[The history machine simulates $M$]
  \label{thm:histTM-simulates}
  \lean{Sagitta.Arrow.TM.Bennett.histTM_simulates, Sagitta.Arrow.TM.Bennett.histTM_halted_iff, Sagitta.Arrow.TM.Bennett.runFrom_clean}
  \leanok
  \uses{def:histTM}
  After $t$ steps from the initial configuration, the projection of the history machine is $M$
  after $t$ steps, and one halts exactly when the other does.
\end{theorem}

\begin{proof}
  \leanok
  Induction on $t$: the history stays clean.
\end{proof}

\begin{theorem}[The history machine is reversible]
  \label{thm:histTM-reversible}
  \lean{Sagitta.Arrow.TM.Bennett.histTM_isReversible, Sagitta.Arrow.TM.Bennett.histTM_isReversible_iff}
  \leanok
  \uses{def:histTM, lem:apply-injective, lem:staysOnBlank}
  If $M$ stays on a blank of its input and the code of the records is injective, the history
  machine is reversible. Conversely, for an injective code, the history machine is reversible only
  if $M$ stays on a blank: the hypothesis is necessary.
\end{theorem}

\begin{proof}
  \leanok
  Two running configurations with the same successor took the same kind of step (the state and the
  last output symbol tell); two good steps wrote the same code, hence the same record, hence the
  same action on the same symbols, which is locally reversible.
\end{proof}

\begin{theorem}[Step 2, packaged]
  \label{thm:history}
  \lean{Sagitta.Arrow.TM.Bennett.history}
  \leanok
  \uses{thm:histTM-simulates, thm:histTM-reversible}
  A machine with finitely many states that stays on a blank of its input has a reversible machine
  with more tapes and finitely many states that projects to it step by step.
\end{theorem}

\begin{proof}
  \leanok
  The history machine with the one-hot code.
\end{proof}

\section{Step 3: the copy}

\begin{definition}[The copying and the emitting machines]
  \label{def:copyTM}
  \lean{Sagitta.Arrow.TM.Bennett.copyTM, Sagitta.Arrow.TM.Bennett.emitTM}
  \leanok
  \uses{def:cslib-tm}
  The copying machine walks a word on a source tape and applies to each cell of a destination tape
  the swap of the blank with the symbol read (the three-valued version of Bennett's exclusive or),
  then walks back. The emitting machine does the same with the output tape as destination.
  \source{\cite{Bennett73}, the copy of the output.}
\end{definition}

\begin{theorem}[The copy is reversible and correct]
  \label{thm:copy}
  \lean{Sagitta.Arrow.TM.Bennett.copyTM_isReversible, Sagitta.Arrow.TM.Bennett.copyTM_self_not_isReversible, Sagitta.Arrow.TM.Bennett.copyTM_correct, Sagitta.Arrow.TM.Bennett.emitTM_isReversible, Sagitta.Arrow.TM.Bennett.emitTM_correct}
  \leanok
  \uses{def:copyTM, def:isReversible}
  The copying machine is reversible when source and destination differ (and not when they are the
  same tape). Started on a word $w$ with blanks around it and $|w|$ blank destination cells, after
  $2|w| + 2$ steps it has halted with $w$ copied and nothing else changed. The emitting machine is
  reversible and appends $w$ to the output in $2|w| + 2$ steps.
\end{theorem}

\begin{proof}
  \leanok
  The successor tells the state and the direction; the untouched source cell tells the symbol, and
  with it the swap is undone.
\end{proof}

\section{Step 4: the inverse}

\begin{definition}[The inverse machines]
  \label{def:invTM}
  \lean{Sagitta.Arrow.TM.Bennett.invTM, Sagitta.Arrow.TM.Bennett.InputFixed, Sagitta.Arrow.TM.Bennett.uncTM}
  \leanok
  \uses{def:histTM}
  The inverse of the history machine reads the last record, undoes the action of $M$ and erases the
  record, moving back. The \emph{uncomputing} machine is a reversible version of it, for machines
  that never move their input head (\emph{input-fixed}).
  \source{\cite{Bennett73}, the uncomputation.}
\end{definition}

\begin{theorem}[The inverse is correct]
  \label{thm:inverse}
  \lean{Sagitta.Arrow.TM.Bennett.invTM_correct, Sagitta.Arrow.TM.Bennett.inverse}
  \leanok
  \uses{def:invTM, thm:histTM-simulates, lem:staysOnBlank}
  If the history machine runs $t$ steps without halting before, its inverse, started on the
  configuration reached, comes back to the initial configuration, halted, in $2t + 2$ steps, with
  the output untouched.
\end{theorem}

\begin{proof}
  \leanok
  Induction on $t$, one record per step.
\end{proof}

\begin{theorem}[The uncomputation is reversible and correct]
  \label{thm:uncompute}
  \lean{Sagitta.Arrow.TM.Bennett.InputFixed.staysOnBlank, Sagitta.Arrow.TM.Bennett.uncTM_isReversible, Sagitta.Arrow.TM.Bennett.uncTM_correct, Sagitta.Arrow.TM.Bennett.uncompute}
  \leanok
  \uses{def:invTM, def:isReversible, thm:histTM-simulates}
  The uncomputing machine is reversible and, for an input-fixed $M$, from the mirror of the
  configuration reached after $t$ running steps it returns to the initial configuration, halted, in
  $3t + 2$ steps.
\end{theorem}

\begin{proof}
  \leanok
  Reversibility by cases on the state; correctness by induction on $t$.
\end{proof}

\section{Step 5: composition}

\begin{definition}[A distinguishable hand-off]
  \label{def:handoff}
  \lean{Sagitta.Arrow.TM.Bennett.HandOff, Sagitta.Arrow.TM.Bennett.NoReentry}
  \leanok
  \uses{def:cslib-plumbing}
  The hand-off from $M_0$ to $M_1$ is \emph{distinguishable} when no running step of $M_1$ reaches a
  configuration where a halting step of $M_0$ hands over to $M_1$. A machine has \emph{no re-entry}
  when no transition enters its initial state.
  \own{}
\end{definition}

\begin{theorem}[Composition keeps reversibility]
  \label{thm:seq-reversible}
  \lean{Sagitta.Arrow.TM.Bennett.seq_isReversible, Sagitta.Arrow.TM.Bennett.seq_not_isReversible, Sagitta.Arrow.TM.Bennett.HandOff.of_noReentry, Sagitta.Arrow.TM.Bennett.seq_runFrom}
  \leanok
  \uses{def:handoff, def:isReversible, def:cslib-plumbing}
  cslib's sequential composition of two reversible machines is reversible when the hand-off is
  distinguishable (no re-entry suffices), and a clash at the hand-off breaks reversibility: the
  condition is necessary. The run of the composition is the run of $M_0$ followed by the run of
  $M_1$.
\end{theorem}

\begin{proof}
  \leanok
  Case analysis on the side of the sum of states.
\end{proof}

\begin{lemma}[The gate before the emission]
  \label{lem:gate}
  \lean{Sagitta.Arrow.TM.Bennett.gateTM, Sagitta.Arrow.TM.Bennett.gateTM_isReversible, Sagitta.Arrow.TM.Bennett.gate_emit_handOff, Sagitta.Arrow.TM.Bennett.uncTM_noReentry}
  \leanok
  \uses{def:handoff, def:copyTM, def:invTM}
  A reversible gate before the emitting machine makes that hand-off distinguishable, and the
  uncomputing machine never re-enters its initial state.
\end{lemma}

\begin{proof}
  \leanok
  After the gate the cell left of the source head is blank; after a forward step of the emission it
  holds a symbol.
\end{proof}

\section{Step 6: Bennett's machine}

\begin{definition}[Bennett's machine]
  \label{def:bennettTM}
  \lean{Sagitta.Arrow.TM.Bennett.emitStage, Sagitta.Arrow.TM.Bennett.bennettTM}
  \leanok
  \uses{def:histTM, def:copyTM, def:invTM, lem:gate}
  Compute with a history, emit the result, uncompute.
  \source{\cite{Bennett73}.}
\end{definition}

\begin{theorem}[Bennett's machine is reversible and correct]
  \label{thm:bennettTM}
  \lean{Sagitta.Arrow.TM.Bennett.bennettTM_isReversible, Sagitta.Arrow.TM.Bennett.bennettTM_correct, Sagitta.Arrow.TM.Bennett.bennettTM_computesInTimeAndSpace}
  \leanok
  \uses{def:bennettTM, thm:histTM-reversible, thm:copy, thm:uncompute, thm:seq-reversible, def:cslib-time-space}
  Bennett's machine is reversible. If an input-fixed $M$ halts at time $t$ with a word $w$ on a work
  tape, Bennett's machine halts by time $4t + 2|w| + 6$ with every work tape blank, every head back
  where it started, and the output of $M$ followed by $w$; it does so in space linear in the time.
\end{theorem}

\begin{proof}
  \leanok
  Compose the three stages and add their times.
\end{proof}

\section{Step 7: preparing the input, the normal form, and the theorem}

\begin{definition}[The preparing machine]
  \label{def:prepTM}
  \lean{Sagitta.Arrow.TM.Bennett.prepTM}
  \leanok
  \uses{def:cslib-tm}
  It copies the input onto a work tape (the virtual input) and leaves a mark on a flag tape, so that
  the rest of the computation never moves the real input head.
  \own{}
\end{definition}

\begin{theorem}[The preparation is reversible and correct]
  \label{thm:prepTM}
  \lean{Sagitta.Arrow.TM.Bennett.prepTM_isReversible, Sagitta.Arrow.TM.Bennett.prepTM_correct}
  \leanok
  \uses{def:prepTM, def:isReversible}
  The preparing machine is reversible and, after $2|x| + 3$ steps, has halted with the input $x$ on
  the virtual input tape, the mark on the flag tape, both heads back at $0$, and nothing else changed.
\end{theorem}

\begin{proof}
  \leanok
  Two phases, forward and back, each an induction on the position.
\end{proof}

\begin{definition}[The normal form]
  \label{def:normTM}
  \lean{Sagitta.Arrow.TM.Bennett.normTM, Sagitta.Arrow.TM.Bennett.InputBlind}
  \leanok
  \uses{def:cslib-plumbing}
  The normal form of $M$ does nothing for one step, then runs $M$ reading its input from the virtual
  input tape (cslib's $\mathrm{inputFromTape}$) and writing its output on a work tape (cslib's
  $\mathrm{outputToTape}$), then rewinds that tape (cslib's $\mathrm{rewindWork}$). It never moves
  nor reads its real input head.
  \own{}
\end{definition}

\begin{theorem}[The normal form is input-fixed and correct]
  \label{thm:normTM}
  \lean{Sagitta.Arrow.TM.Bennett.normTM_inputFixed, Sagitta.Arrow.TM.Bennett.normTM_inputBlind, Sagitta.Arrow.TM.Bennett.normTM_noReentry, Sagitta.Arrow.TM.Bennett.normTM_run}
  \leanok
  \uses{def:normTM, def:invTM, def:handoff}
  The normal form is input-fixed, input-blind and has no re-entry. If $M$ has halted by time $t$ with
  output $y$, the normal form, started on the prepared tapes, has halted by time $1 + t + |y| + 2$
  with $y$ on its output work tape from cell $0$, its head at $0$, and nothing on the real output.
\end{theorem}

\begin{proof}
  \leanok
  cslib's lemmas on the plumbing, plus a simulation lemma for the moved input head.
\end{proof}

\begin{definition}[The reversible machine for $M$]
  \label{def:revTM}
  \lean{Sagitta.Arrow.TM.Bennett.revTM}
  \leanok
  \uses{def:prepTM, def:normTM, def:bennettTM}
  Prepare the virtual input, then run Bennett's machine on the normal form of $M$.
  \own{}
\end{definition}

\begin{theorem}[The reversible machine for $M$ is reversible and correct]
  \label{thm:revTM}
  \lean{Sagitta.Arrow.TM.Bennett.revTM_isReversible, Sagitta.Arrow.TM.Bennett.revTM_correct}
  \leanok
  \uses{def:revTM, thm:prepTM, thm:normTM, thm:bennettTM, thm:seq-reversible}
  It is reversible, and if $M$ halts on $x$ by time $t$ with output $y$, it halts with output $y$ by
  time $2|x| + 3 + 4(1 + t + |y| + 2) + 2|y| + 6$.
\end{theorem}

\begin{proof}
  \leanok
  Compose the preparation with Bennett's machine on the normal form.
\end{proof}

\begin{theorem}[Bennett's theorem (1973), for cslib's machines]
  \label{thm:bennett-holds}
  \lean{Sagitta.Arrow.TM.bennett_holds}
  \leanok
  \uses{def:bennett, thm:revTM, def:polytime}
  Bennett's statement holds: every function computable in polynomial time is computed by a
  reversible machine with finitely many states in time $2n + 4t + 6|y| + 21 \le 2n + 10\,t + 21$ (for a
  machine running in time $t$ with output $y$ on inputs of length $n$) and space linear in that
  time.
  \source{\cite{Bennett73}.}
  \fidelity{clean tapes being added to the statement; see Definition~\ref{def:bennett} and the
  Fidelity notes of the README. The construction itself uncomputes and leaves every work tape blank
  (Theorem~\ref{thm:bennettTM}); the $4t$ matches Bennett's count.}
\end{theorem}

\begin{proof}
  \leanok
  \uses{thm:revTM}
  Apply Theorem~\ref{thm:revTM} to a machine that computes the function, at its time bound.
\end{proof}

\section{Bennett and the arrow}

\begin{theorem}[Every $\cNP$ verifier runs on a reversible machine]
  \label{thm:reversible-verifier}
  \lean{Sagitta.Arrow.TM.reversible_verifier, Sagitta.Arrow.TM.reversible_verifier_holds}
  \leanok
  \uses{thm:bennett-holds, def:NP}
  For every problem in $\cNP$ there is a verifier computed in polynomial time and space by a
  reversible machine: the dynamics of checking a witness carries no arrow.
  \own{}
\end{theorem}

\begin{proof}
  \leanok
  Apply Bennett's theorem to a polynomial-time verifier.
\end{proof}

\begin{theorem}[HD and Bennett]
  \label{thm:decision-bennett}
  \lean{Sagitta.Arrow.TM.decision_bennett, Sagitta.Arrow.TM.decision_bennett_holds}
  \leanok
  \uses{thm:reversible-verifier, thm:hd-iff-arrow-tm}
  From HD alone: the arrow holds, and every $\cNP$ verifier runs on a reversible machine in
  polynomial time and space. As the past hypothesis and reversible dynamics give the second law of
  thermodynamics, HD and Bennett give the arrow and locate it: it is in the reduction to one bit, not
  in the dynamics.
  \own{}
\end{theorem}

\begin{proof}
  \leanok
  The conjunction of Theorems~\ref{thm:hd-iff-arrow-tm} and~\ref{thm:reversible-verifier}.
\end{proof}

\chapter{The consequences}
\label{chap:consequences}

Every consequence takes its hypothesis as an explicit argument: the arrow, HD (equivalent to it)
or law 2.2 (stronger). None of the hypotheses is proved, and none is an axiom.

\section{In every model}

\begin{theorem}[Consequences of the arrow, in every model]
  \label{thm:model-consequences}
  \lean{Sagitta.Arrow.Model.not_InP_of_NPComplete, Sagitta.Arrow.Model.not_searchEasy, Sagitta.Arrow.Model.not_reduces_of_InP}
  \leanok
  \uses{def:arrow, thm:kelvin, lem:npcomplete-reduces}
  Assume the arrow. Then no $\cNP$-complete problem is in $\cP$, search is not easy, and no
  $\cNP$-complete problem reduces to a problem in $\cP$.
\end{theorem}

\begin{proof}
  \leanok
  An $\cNP$-complete problem in $\cP$ gives $\cP = \cNP$ (Lemma~\ref{lem:npcomplete-reduces}).
\end{proof}

\begin{theorem}[Consequences of HD, in every model]
  \label{thm:model-consequences-hd}
  \lean{Sagitta.Arrow.Model.not_InP_of_NPComplete_of_decision, Sagitta.Arrow.Model.not_searchEasy_of_decision, Sagitta.Arrow.Model.not_reduces_of_InP_of_decision}
  \leanok
  \uses{def:hd, thm:hd-iff-arrow, thm:model-consequences}
  The same consequences follow from HD.
\end{theorem}

\begin{proof}
  \leanok
  HD implies the arrow.
\end{proof}

\begin{proposition}[Control: without the arrow the consequences fail]
  \label{prop:consequences-need-arrow}
  \lean{Sagitta.Arrow.consequences_need_arrow}
  \leanok
  \uses{def:two-models, def:arrow}
  In the collapsed model the only problem is $\cNP$-complete and in $\cP$, and every search problem
  is easy.
\end{proposition}

\begin{proof}
  \leanok
  Compute.
\end{proof}

\section{SAT and 3-SAT}

\begin{theorem}[$\SAT$ and $\TSAT$ are not in $\cP$]
  \label{thm:sat-notmem-p}
  \lean{Sagitta.Arrow.TM.SAT_notMem_P, Sagitta.Arrow.TM.ThreeSAT_notMem_P}
  \leanok
  \uses{def:arrow-tm, thm:arrow-iff-tm, thm:cook-levin}
  Assume the arrow. Then $\SAT \notin \cP$ and $\TSAT \notin \cP$.
  \source{the Cook--Levin theorem \cite{Cook71, Levin73} and the reduction of $\SAT$ to $\TSAT$
  \cite{Karp72}, both proved in cslib-frontier.}
\end{theorem}

\begin{proof}
  \leanok
  \uses{thm:cook-levin}
  $\SAT$ is $\cNP$-complete and reduces to $\TSAT$.
\end{proof}

\begin{theorem}[From HD]
  \label{thm:sat-notmem-p-hd}
  \lean{Sagitta.Arrow.TM.SAT_notMem_P_of_decision, Sagitta.Arrow.TM.ThreeSAT_notMem_P_of_decision, Sagitta.Arrow.TM.not_searchEasy_of_decision}
  \leanok
  \uses{def:hd-tm, thm:hd-iff-arrow-tm, thm:sat-notmem-p}
  Assume HD. Then $\SAT, \TSAT \notin \cP$, and some $\cFNP$ search problem is not in $\cFP$.
\end{theorem}

\begin{proof}
  \leanok
  HD implies the arrow.
\end{proof}

\section{An $\cNP$-intermediate problem}

\begin{theorem}[Ladner, under the arrow]
  \label{thm:intermediate}
  \lean{Sagitta.Arrow.TM.exists_intermediate, Sagitta.Arrow.TM.exists_intermediate_model}
  \leanok
  \uses{def:arrow-tm, thm:arrow-iff-tm, thm:ladner, lem:npComplete-iff}
  Assume the arrow. Then some problem over binary words is in $\cNP$, not in $\cP$ and not
  $\cNP$-complete; in the vocabulary of the model, the Turing-machine model has an intermediate
  problem.
  \source{\cite{Ladner75}; the proof is cslib-frontier's (Theorem~\ref{thm:ladner}).}
\end{theorem}

\begin{proof}
  \leanok
  \uses{thm:ladner}
  The arrow gives an encoding with $\cP \neq \cNP$; apply Ladner's theorem.
\end{proof}

\begin{theorem}[From HD]
  \label{thm:intermediate-hd}
  \lean{Sagitta.Arrow.TM.exists_intermediate_of_decision, Sagitta.Arrow.TM.exists_intermediate_model_of_decision}
  \leanok
  \uses{def:hd-tm, thm:hd-iff-arrow-tm, thm:intermediate}
  Assume HD. Then there is an $\cNP$-intermediate problem.
\end{theorem}

\begin{proof}
  \leanok
  HD implies the arrow.
\end{proof}

\section{The Ising ground state}

\begin{definition}[$\ISING$]
  \label{def:ising}
  \lean{Sagitta.Arrow.TM.Ising.Coupling, Sagitta.Arrow.TM.Ising.IsingInstance, Sagitta.Arrow.TM.Ising.energy, Sagitta.Arrow.TM.Ising.ISING}
  \leanok
  An instance is a list of couplings $J_e$ between sites $u_e, v_e \in \mathbb{N}$ and a threshold
  $t \in \mathbb{Z}$. The energy of a configuration $\sigma : \mathbb{N} \to \{\pm 1\}$ is
  $H_J(\sigma) = -\sum_e J_e\, \sigma_{u_e} \sigma_{v_e}$ (no external field). $\ISING$ is the set of
  instances with some configuration of energy at most $t$.
  \source{\cite{Barahona82}.}
\end{definition}

\begin{theorem}[$\ISING$ is $\cNP$-complete]
  \label{thm:ising-npcomplete}
  \lean{Sagitta.Arrow.TM.Ising.ISING_mem_NP, Sagitta.Arrow.TM.Ising.maxCut_iff_ising, Sagitta.Arrow.TM.Ising.polyTimeReducible_MaxCut_ISING, Sagitta.Arrow.TM.Ising.ISING_npComplete, Sagitta.Arrow.TM.Ising.P_eq_NP_iff_ISING_mem_P, Sagitta.Arrow.TM.Ising.ISING_mem_P_iff_SAT_mem_P}
  \leanok
  \uses{def:ising, thm:maxcut, def:reduction, thm:cook-levin}
  $\ISING \in \cNP$; $\MAXCUT \red \ISING$ through Barahona's identity
  $H_J(\sigma) = 2\,\mathrm{cut}_J(\sigma) - \sum_e J_e$; hence $\ISING$ is $\cNP$-complete,
  $\cP = \cNP$ if and only if $\ISING \in \cP$, and $\ISING \in \cP$ if and only if $\SAT \in \cP$.
  \source{\cite{Barahona82}.}
  \fidelity{this is $\cNP$-completeness on arbitrary graphs, with integer couplings written in
  binary and no field, by the reduction from $\MAXCUT$. Barahona proves the much stronger result on
  three-dimensional lattices with couplings in $\{-1, 0, 1\}$, which is not formalized here.}
\end{theorem}

\begin{proof}
  \leanok
  \uses{thm:maxcut}
  A verifier checks the energy of the configuration given by the endpoints in the certificate; the
  reduction turns every edge of weight $w$ into the antiferromagnetic coupling $J = -w$, with
  threshold $W - 2k$ ($W$ the total weight, $k$ the cut bound).
\end{proof}

\begin{theorem}[$\ISING$ is not in $\cP$]
  \label{thm:ising-notmem-p}
  \lean{Sagitta.Arrow.TM.ISING_notMem_P, Sagitta.Arrow.TM.ISING_notMem_P_of_decision}
  \leanok
  \uses{thm:ising-npcomplete, thm:sat-notmem-p, def:hd-tm, thm:hd-iff-arrow-tm}
  Under the arrow, and so under HD, $\ISING \notin \cP$.
\end{theorem}

\begin{proof}
  \leanok
  $\ISING \in \cP$ would put $\SAT$ in $\cP$.
\end{proof}

\section{Under law 2.2}

\begin{theorem}[Consequences of law 2.2]
  \label{thm:law22-consequences}
  \lean{Sagitta.Arrow.TM.ThreeSAT_notMem_TIME, Sagitta.Arrow.TM.arrow_of_law22, Sagitta.Arrow.TM.decisionHypothesis_of_law22}
  \leanok
  \uses{def:law22, thm:law22-arrow, thm:hd-iff-arrow-tm}
  Assume law 2.2. Then for every $\varepsilon > 0$, $\TSAT$ is not decided within
  $C \cdot \lceil 2^{n^{1-\varepsilon}} \rceil + C$ steps for any constant $C$; the arrow holds; and
  HD holds. All the consequences above follow.
\end{theorem}

\begin{proof}
  \leanok
  The first statement is the definition, unfolded; the others are Theorem~\ref{thm:law22-arrow} and
  the equivalence of HD with the arrow.
\end{proof}

\begin{thebibliography}{99}

\bibitem{AroraBarak}
S.~Arora and B.~Barak.
\newblock \emph{Computational Complexity: A Modern Approach}.
\newblock Cambridge University Press, 2009.

\bibitem{Barahona82}
F.~Barahona.
\newblock On the computational complexity of Ising spin glass models.
\newblock \emph{Journal of Physics A: Mathematical and General}, 15(10):3241--3253, 1982.

\bibitem{Bennett73}
C.~H. Bennett.
\newblock Logical reversibility of computation.
\newblock \emph{IBM Journal of Research and Development}, 17(6):525--532, 1973.

\bibitem{Cook71}
S.~A. Cook.
\newblock The complexity of theorem-proving procedures.
\newblock In \emph{Proceedings of the Third Annual ACM Symposium on Theory of Computing (STOC)},
pages 151--158, 1971.

\bibitem{ETHdead}
MIT Hardness Group, J.~Brunner, E.~D. Demaine, F.~Frei, T.~Gomez, M.~Hecher, J.~Li, J.~Lynch, et al.
\newblock ETH is dead, long live SH.
\newblock arXiv:2611.02650, November 2026.

\bibitem{Goldreich}
O.~Goldreich.
\newblock \emph{Computational Complexity: A Conceptual Perspective}.
\newblock Cambridge University Press, 2008.

\bibitem{Karp72}
R.~M. Karp.
\newblock Reducibility among combinatorial problems.
\newblock In R.~E. Miller and J.~W. Thatcher, editors, \emph{Complexity of Computer
Computations}, pages 85--103. Plenum Press, 1972.

\bibitem{Koopman31}
B.~O. Koopman.
\newblock Hamiltonian systems and transformation in Hilbert space.
\newblock \emph{Proceedings of the National Academy of Sciences}, 17(5):315--318, 1931.

\bibitem{Ladner75}
R.~E. Ladner.
\newblock On the structure of polynomial time reducibility.
\newblock \emph{Journal of the ACM}, 22(1):155--171, 1975.

\bibitem{Levin73}
L.~A. Levin.
\newblock Universal sequential search problems.
\newblock \emph{Problems of Information Transmission}, 9(3):265--266, 1973.

\bibitem{Papadimitriou}
C.~H. Papadimitriou.
\newblock \emph{Computational Complexity}.
\newblock Addison-Wesley, 1994.

% The references of the chapters of the map, one file per chapter.
% References of chapters/map/machines.tex: \bibitem entries, in alphabetical order of key.
% An entry that has not been checked against the publication ends its first line with: % unverified

\bibitem{BookGreibach70} % unverified
R.~V. Book and S.~A. Greibach.
\newblock Quasi-realtime languages.
\newblock \emph{Mathematical Systems Theory}, 4(2):97--111, 1970.

\bibitem{ChandraKozenStockmeyer81} % unverified
A.~K. Chandra, D.~C. Kozen, and L.~J. Stockmeyer.
\newblock Alternation.
\newblock \emph{Journal of the ACM}, 28(1):114--133, 1981.

\bibitem{Cook06} % unverified
S.~A. Cook.
\newblock The P versus NP problem.
\newblock In J.~Carlson, A.~Jaffe, and A.~Wiles, editors, \emph{The Millennium Prize Problems},
pages 87--104. Clay Mathematics Institute and American Mathematical Society, 2006.

\bibitem{KarpLipton80} % unverified
R.~M. Karp and R.~J. Lipton.
\newblock Some connections between nonuniform and uniform complexity classes.
\newblock In \emph{Proceedings of the 12th Annual ACM Symposium on Theory of Computing (STOC)},
pages 302--309, 1980.

\bibitem{MeyerStockmeyer72} % unverified
A.~R. Meyer and L.~J. Stockmeyer.
\newblock The equivalence problem for regular expressions with squaring requires exponential space.
\newblock In \emph{Proceedings of the 13th Annual IEEE Symposium on Switching and Automata Theory},
pages 125--129, 1972.

\bibitem{Savitch70} % unverified
W.~J. Savitch.
\newblock Relationships between nondeterministic and deterministic tape complexities.
\newblock \emph{Journal of Computer and System Sciences}, 4(2):177--192, 1970.

\bibitem{SeiferasFischerMeyer78} % unverified
J.~I. Seiferas, M.~J. Fischer, and A.~R. Meyer.
\newblock Separating nondeterministic time complexity classes.
\newblock \emph{Journal of the ACM}, 25(1):146--167, 1978.

\bibitem{Shannon49} % unverified
C.~E. Shannon.
\newblock The synthesis of two-terminal switching circuits.
\newblock \emph{Bell System Technical Journal}, 28(1):59--98, 1949.

\bibitem{Zak83} % unverified
S.~\v{Z}\'ak.
\newblock A Turing machine time hierarchy.
\newblock \emph{Theoretical Computer Science}, 26(3):327--333, 1983.

% References of chapters/map/search.tex: \bibitem entries, in alphabetical order of key.
% An entry that has not been checked against the publication ends its first line with: % unverified

\bibitem{BeckmannBussPollett03} % unverified
A.~Beckmann, S.~R. Buss, and C.~Pollett.
\newblock Ordinal notations and well-orderings in bounded arithmetic.
\newblock \emph{Annals of Pure and Applied Logic}, 120:197--223, 2003.

\bibitem{FortnowVanMelkebeek00} % unverified
L.~Fortnow and D.~van Melkebeek.
\newblock Time-space tradeoffs for nondeterministic computation.
\newblock In \emph{Proceedings of the 15th Annual IEEE Conference on Computational Complexity},
2000.

\bibitem{JohnsonPapadimitriouYannakakis88} % unverified
D.~S. Johnson, C.~H. Papadimitriou, and M.~Yannakakis.
\newblock How easy is local search?
\newblock \emph{Journal of Computer and System Sciences}, 37(1):79--100, 1988.

\bibitem{LiptonViglas99} % unverified
R.~J. Lipton and A.~Viglas.
\newblock On the complexity of {SAT}.
\newblock In \emph{Proceedings of the 40th Annual Symposium on Foundations of Computer Science
(FOCS)}, 1999.

\bibitem{MegiddoPapadimitriou91} % unverified
N.~Megiddo and C.~H. Papadimitriou.
\newblock On total functions, existence theorems and computational complexity.
\newblock \emph{Theoretical Computer Science}, 81(2):317--324, 1991.

\bibitem{MoserScheder11} % unverified
R.~A. Moser and D.~Scheder.
\newblock A full derandomization of Sch{\"o}ning's $k$-{SAT} algorithm.
\newblock In \emph{Proceedings of the 43rd Annual ACM Symposium on Theory of Computing (STOC)},
2011.

\bibitem{Schoening99} % unverified
U.~Sch{\"o}ning.
\newblock A probabilistic algorithm for $k$-{SAT} and constraint satisfaction problems.
\newblock In \emph{Proceedings of the 40th Annual Symposium on Foundations of Computer Science
(FOCS)}, 1999.

% References of chapters/map/dynamics.tex: \bibitem entries, in alphabetical order of key.
% An entry that has not been checked against the publication ends its first line with: % unverified

% References of chapters/map/information.tex: \bibitem entries, in alphabetical order of key.
% An entry that has not been checked against the publication ends its first line with: % unverified

\bibitem{Antunes06} % unverified
L.~Antunes, L.~Fortnow, D.~van Melkebeek and N.~V. Vinodchandran.
\newblock Computational depth: concept and applications.
\newblock \emph{Theoretical Computer Science}, 354(3):391--404, 2006.

\bibitem{Bennett88} % unverified
C.~H. Bennett.
\newblock Logical depth and physical complexity.
\newblock In R.~Herken, editor, \emph{The Universal Turing Machine: A Half-Century Survey},
pages 227--257. Oxford University Press, 1988.

\bibitem{Blum67} % unverified
M.~Blum.
\newblock A machine-independent theory of the complexity of recursive functions.
\newblock \emph{Journal of the ACM}, 14(2):322--336, 1967.

\bibitem{Kolmogorov65} % unverified
A.~N. Kolmogorov.
\newblock Three approaches to the quantitative definition of information.
\newblock \emph{Problems of Information Transmission}, 1(1):1--7, 1965.

\bibitem{Valiant79} % unverified
L.~G. Valiant.
\newblock The complexity of computing the permanent.
\newblock \emph{Theoretical Computer Science}, 8(2):189--201, 1979.

% References of chapters/map/arithmetic.tex: \bibitem entries, in alphabetical order of key.
% An entry that has not been checked against the publication ends its first line with: % unverified

\bibitem{BKO20} % unverified
J.~Bydzovsk\'y, J.~Kraj\'{\i}\v{c}ek and I.~C. Oliveira.
\newblock Consistency of circuit lower bounds with bounded theories.
\newblock \emph{Logical Methods in Computer Science}, 16(2), 2020.

\bibitem{Buss86} % unverified
S.~R. Buss.
\newblock \emph{Bounded Arithmetic}.
\newblock Bibliopolis, Naples, 1986.

\bibitem{Buss98} % unverified
S.~R. Buss.
\newblock First-order proof theory of arithmetic.
\newblock In S.~R. Buss, editor, \emph{Handbook of Proof Theory}. Elsevier, 1998.

\bibitem{Cobham65} % unverified
A.~Cobham.
\newblock The intrinsic computational difficulty of functions.
\newblock In \emph{Logic, Methodology and Philosophy of Science}. North-Holland, 1965.

\bibitem{Cook75} % unverified
S.~A. Cook.
\newblock Feasibly constructive proofs and the propositional calculus.
\newblock In \emph{Proceedings of the Seventh Annual ACM Symposium on Theory of Computing}, 1975.

\bibitem{CookReckhow79} % unverified
S.~A. Cook and R.~A. Reckhow.
\newblock The relative efficiency of propositional proof systems.
\newblock \emph{Journal of Symbolic Logic}, 44(1):36--50, 1979.

\bibitem{Godel30} % unverified
K.~G\"odel.
\newblock Die Vollst\"andigkeit der Axiome des logischen Funktionenkalk\"uls.
\newblock \emph{Monatshefte f\"ur Mathematik und Physik}, 37, 1930.

\bibitem{Godel31} % unverified
K.~G\"odel.
\newblock \"Uber formal unentscheidbare S\"atze der Principia Mathematica und verwandter Systeme I.
\newblock \emph{Monatshefte f\"ur Mathematik und Physik}, 38, 1931.

\bibitem{Henkin49} % unverified
L.~Henkin.
\newblock The completeness of the first-order functional calculus.
\newblock \emph{Journal of Symbolic Logic}, 14(3), 1949.

\bibitem{KPT91} % unverified
J.~Kraj\'{\i}\v{c}ek, P.~Pudl\'ak and G.~Takeuti.
\newblock Bounded arithmetic and the polynomial hierarchy.
\newblock \emph{Annals of Pure and Applied Logic}, 52(1--2):143--153, 1991.

\bibitem{Lob55} % unverified
M.~H. L\"ob.
\newblock Solution of a problem of Leon Henkin.
\newblock \emph{Journal of Symbolic Logic}, 20(2), 1955.

\bibitem{Parikh71} % unverified
R.~J. Parikh.
\newblock Existence and feasibility in arithmetic.
\newblock \emph{Journal of Symbolic Logic}, 36(3), 1971.

\bibitem{ParisWilkie85} % unverified
J.~B. Paris and A.~J. Wilkie.
\newblock Counting problems in bounded arithmetic.
\newblock In \emph{Methods in Mathematical Logic}, Lecture Notes in Mathematics 1130. Springer, 1985.

\bibitem{Razborov95} % unverified
A.~A. Razborov.
\newblock Unprovability of lower bounds on circuit size in certain fragments of bounded arithmetic.
\newblock \emph{Izvestiya: Mathematics}, 59(1), 1995.

\bibitem{Schaefer78} % unverified
T.~J. Schaefer.
\newblock The complexity of satisfiability problems.
\newblock In \emph{Proceedings of the Tenth Annual ACM Symposium on Theory of Computing (STOC)}, 1978.

\bibitem{Jeavons98} % unverified
P.~Jeavons.
\newblock On the algebraic structure of combinatorial problems.
\newblock \emph{Theoretical Computer Science}, 200(1--2), 1998.

\bibitem{Geiger68} % unverified
D.~Geiger.
\newblock Closed systems of functions and predicates.
\newblock \emph{Pacific Journal of Mathematics}, 27(1), 1968.

\bibitem{BodnarchukEtAl69} % unverified
V.~G. Bodnarchuk, L.~A. Kaluzhnin, V.~N. Kotov and B.~A. Romov.
\newblock Galois theory for Post algebras.
\newblock \emph{Kibernetika}, 1969.

\bibitem{DowlingGallier84} % unverified
W.~F. Dowling and J.~H. Gallier.
\newblock Linear-time algorithms for testing the satisfiability of propositional Horn formulae.
\newblock \emph{Journal of Logic Programming}, 1(3), 1984.

\bibitem{GallantMaierStorer80} % unverified
J.~Gallant, D.~Maier and J.~A. Storer.
\newblock On finding minimal length superstrings.
\newblock \emph{Journal of Computer and System Sciences}, 20(1), 1980.

\bibitem{Pratt75} % unverified
V.~R. Pratt.
\newblock Every prime has a succinct certificate.
\newblock \emph{SIAM Journal on Computing}, 4(3):214--220, 1975.

\bibitem{BeyersdorffGalesiLauria13} % unverified
O.~Beyersdorff, N.~Galesi and M.~Lauria.
\newblock A lower bound for the pigeonhole principle in tree-like resolution by asymmetric Prover-Delayer games.
\newblock \emph{Information Processing Letters}, 110(23), 2010.

\bibitem{Pudlak97} % unverified
P.~Pudl\'ak.
\newblock Lower bounds for resolution and cutting plane proofs and monotone computations.
\newblock \emph{Journal of Symbolic Logic}, 62(3), 1997.

\bibitem{Krajicek97} % unverified
J.~Kraj\'i\v{c}ek.
\newblock Interpolation theorems, lower bounds for proof systems, and independence results for bounded arithmetic.
\newblock \emph{Journal of Symbolic Logic}, 62(2), 1997.

\bibitem{AlonBoppana87} % unverified
N.~Alon and R.~B. Boppana.
\newblock The monotone circuit complexity of Boolean functions.
\newblock \emph{Combinatorica}, 7(1), 1987.

\bibitem{BenDavidHalevi91} % unverified
S.~Ben-David and S.~Halevi.
\newblock On the independence of {P} versus {NP}.
\newblock Technical Report 699, Technion, 1991.

\bibitem{Schoning82} % unverified
U.~Sch\"oning.
\newblock A uniform approach to obtain diagonal sets in complexity classes.
\newblock \emph{Theoretical Computer Science}, 18(1), 1982.

% References of chapters/map/barriers.tex: \bibitem entries, in alphabetical order of key.
% An entry that has not been checked against the publication ends its first line with: % unverified


\bibitem{Alon99} % unverified
N.~Alon.
\newblock Combinatorial Nullstellensatz.
\newblock \emph{Combinatorics, Probability and Computing}, 8(1--2):7--29, 1999.

\bibitem{BGS75} % unverified
T.~Baker, J.~Gill and R.~Solovay.
\newblock Relativizations of the $\mathcal{P} =? \mathcal{NP}$ question.
\newblock \emph{SIAM Journal on Computing}, 4(4):431--442, 1975.

\bibitem{Burnside} % unverified
W.~Burnside.
\newblock \emph{Theory of Groups of Finite Order}.
\newblock Cambridge University Press.

\bibitem{Lawvere69} % unverified
F.~W. Lawvere.
\newblock Diagonal arguments and cartesian closed categories.
\newblock In \emph{Category Theory, Homology Theory and their Applications II}, Lecture Notes in Mathematics 92, Springer, 1969.

\bibitem{Turing36} % unverified
A.~M. Turing.
\newblock On computable numbers, with an application to the Entscheidungsproblem.
\newblock \emph{Proceedings of the London Mathematical Society}, s2-42(1):230--265, 1936.


% References of the chapter Physics: the Ising model. Keys of references.tex are not repeated here.

\bibitem{FrancoPaull83} % unverified
J.~Franco and M.~Paull.
\newblock Probabilistic analysis of the Davis Putnam procedure for solving the satisfiability problem.
\newblock \emph{Discrete Applied Mathematics}, 5(1):77--87, 1983.

\bibitem{MezardParisiZecchina02} % unverified
M.~M\'ezard, G.~Parisi and R.~Zecchina.
\newblock Analytic and algorithmic solution of random satisfiability problems.
\newblock \emph{Science}, 297(5582):812--815, 2002.

\bibitem{PaleyZygmund32} % unverified
R.~E. A.~C. Paley and A.~Zygmund.
\newblock On some series of functions (3).
\newblock \emph{Mathematical Proceedings of the Cambridge Philosophical Society}, 28, 1932.

\bibitem{Poincare1890} % unverified
H.~Poincar\'e.
\newblock Sur le probl\`eme des trois corps et les \'equations de la dynamique.
\newblock \emph{Acta Mathematica}, 13:1--270, 1890.

% References of chapters/map/infinite.tex: \bibitem entries, in alphabetical order of key.
% An entry that has not been checked against the publication ends its first line with: % unverified

\bibitem{Goodstein44} % unverified
R.~L. Goodstein.
\newblock On the restricted ordinal theorem.
\newblock \emph{The Journal of Symbolic Logic}, 9(2):33--41, 1944.

\bibitem{Kolmogorov33} % unverified
A.~N. Kolmogorov.
\newblock \emph{Grundbegriffe der Wahrscheinlichkeitsrechnung}.
\newblock Springer, 1933.

\bibitem{LobWainer70} % unverified
M.~H. L\"ob and S.~S. Wainer.
\newblock Hierarchies of number-theoretic functions. I.
\newblock \emph{Archiv f\"ur mathematische Logik und Grundlagenforschung}, 13:39--51, 1970.

\bibitem{Oxtoby80} % unverified
J.~C. Oxtoby.
\newblock \emph{Measure and Category}.
\newblock Graduate Texts in Mathematics 2. Springer, second edition, 1980.

% References of chapters/map/scales.tex: \bibitem entries, in alphabetical order of key.
% An entry that has not been checked against the publication ends its first line with: % unverified

% References of chapters/map/arrow.tex: \bibitem entries, in alphabetical order of key.
% An entry that has not been checked against the publication ends its first line with: % unverified


\end{thebibliography}

